\documentclass{article}
\usepackage[paperwidth=140cm,paperheight=90cm,margin=0cm]{geometry}
\usepackage[T1]{fontenc}
\usepackage{cmbright}
\usepackage{amsmath,amssymb,graphicx,tikz,ragged2e}
\usepackage[hidelinks]{hyperref}
\usepackage{qrcode,multicol}
\usepackage[numbers,sort&compress]{natbib}
\bibliographystyle{posterunsrtnat}
\renewcommand{\bibsection}{}
\renewcommand{\bibfont}{\fontsize{18}{21}\selectfont}
\setlength{\bibsep}{0.12cm}
\definecolor{WhartonBlue}{HTML}{011F5B}
\definecolor{WhartonRed}{HTML}{990000}
\definecolor{Ink}{HTML}{2D2C41}
\definecolor{PaleBlue}{HTML}{EEF3F9}
\definecolor{Rule}{HTML}{D7DFEA}
\pagestyle{empty}
\setlength{\parindent}{0pt}
\newcommand{\bodyfont}{\fontsize{30}{37}\selectfont}
\newcommand{\captionfont}{\fontsize{26}{32}\selectfont\color{Ink}}
\newcommand{\icon}[2][1.3cm]{\raisebox{-0.12em}{\includegraphics[height=#1]{assets/lucide/#2.pdf}}}
\newcommand{\heading}[1]{{\color{WhartonBlue}\fontsize{46}{52}\selectfont\bfseries #1}}
\newcommand{\subhead}[1]{{\color{WhartonBlue}\fontsize{36}{43}\selectfont\bfseries #1}}
% Coordinates in centimetres from the upper-left corner.
\newcommand{\block}[4]{\node[anchor=north west,inner sep=0pt,text width=#3cm] at (#1,-#2) {\RaggedRight\bodyfont\color{Ink}#4};}
\newcommand{\fig}[4]{\node[anchor=north west,inner sep=0pt] at (#1,-#2) {\includegraphics[width=#3cm]{#4}};}
\begin{document}\null
\begin{tikzpicture}[remember picture,overlay,x=1cm,y=1cm]
\begin{scope}[shift={(current page.north west)}]
\fill[WhartonBlue] (0,0) rectangle (140,-0.65);
\fig{102}{2.2}{34}{assets/wharton-statistics.pdf}
\block{4}{2}{106}{{\fontsize{82}{92}\selectfont\bfseries\color{WhartonBlue}Estimating Implicit Regularization in Deep Learning}}
\block{4}{6.1}{132}{{\fontsize{34}{42}\selectfont Joseph H. Rudoler\enspace $\cdot$\enspace Kevin Tan\enspace $\cdot$\enspace Giles Hooker\enspace $\cdot$\enspace Konrad P. Kording\quad {\color{WhartonBlue}\textbar}\quad {\fontsize{28}{35}\selectfont\color{WhartonBlue}University of Pennsylvania}}}
\draw[WhartonBlue,line width=1.6pt] (4,-9.5) -- (136,-9.5);
\draw[Rule,line width=1pt] (47.5,-12) -- (47.5,-83.5);
\draw[Rule,line width=1pt] (92.5,-12) -- (92.5,-83.5);
\fill[PaleBlue,rounded corners=8pt] (4,-12) rectangle (46,-33.2);
\fill[WhartonRed] (4,-12) rectangle (4.3,-33.2);
\block{5.3}{13}{39.4}{\subhead{What regularization is implicit in your training procedure?}\\[0.6cm]
{\bodyfont
\textbf{Why regularization?} Explicit regularization adds penalties to the objective function that steer learning toward certain solutions---for example, those with smaller weights---in ways that can reduce overfitting and improve generalization. Training choices such as early stopping, minibatching, and dropout also bias which solutions are learned, but we often don't know exactly how~\cite{bartlettDeepLearningStatistical2021a,smithOriginImplicitRegularization2021}.\\[0.6cm]
\textbf{Characterizing implicit regularization.} Theory characterizes some of these implicit effects, mostly in simplified settings~\cite{soudryImplicitBiasGradient2018,aroraImplicitRegularizationDeep2019}. We propose an empirical framework for testing theoretical predictions and extending the study of implicit regularization to networks that are not analytically tractable.\\[0.6cm]
\textbf{Our approach.} Work backwards: observe the learned weights and training dynamics, and estimate a parameterized regularizer that explains them. We match loss gradients or update residuals to the gradients of a proposed regularizer and solve for its parameters. \par}}
\block{4}{34}{42}{\heading{\icon{scan-search}\quad Gradient matching}}
\block{4}{36.7}{42}{{\bodyfont
Training minimizes empirical risk (ERM). We model the learned weights as a regularized ERM (RERM) solution, choosing a family $\mathcal R(\theta,D;\Lambda)$ and estimating its parameters $\Lambda$:
\[\hat\theta_{\mathrm{ERM}}=\arg\min_\theta\mathcal L(\theta,D)
\qquad\longrightarrow\qquad
\hat\theta_{\mathrm{RERM}}=\arg\min_\theta\big\{\mathcal L(\theta,D)+\mathcal R(\theta,D;\Lambda)\big\}.\]\par}}
\fig{13}{42.2}{24}{figures/tradeoff-vis.pdf}
\block{4}{58}{42}{{\captionfont An $\ell_2$ penalty shifts the solution, where the empirical loss gradient and regularizer gradient cancel.\par}}
\block{4}{59.8}{42}{\subhead{Estimating from learned weights}\\[0.3cm]
{\bodyfont
\setlength{\abovedisplayskip}{0.25cm}\setlength{\belowdisplayskip}{0.25cm}
At a hypothetical optimum $\hat\theta$ of $\mathcal L+\mathcal R$, the gradients sum to zero:
\[\color{WhartonBlue}-\nabla_\theta\mathcal L(\hat\theta,D)=\nabla_\theta\mathcal R(\hat\theta,D;\Lambda).\]
These $p=\dim(\theta)$ equations let us solve for, or estimate, $\Lambda$:
\[\hat\Lambda=\arg\min_\Lambda\frac{1}{p}\left\|\nabla_\theta\mathcal L+\nabla_\theta\mathcal R_\Lambda\right\|_2^2.\]
\textbf{Example:} For $\mathcal L=\|y-X\theta\|^2/n$ and $\mathcal R=\lambda\|\theta\|^2$,
\[X^\top(y-X\hat\theta)/n=\lambda\hat\theta.\]\par}}
\block{4}{72}{42}{\subhead{Estimating from weight updates}}
\fig{36}{70}{10}{figures/sgd-vs-full-batch.pdf}
\block{4}{74.5}{29.5}{{\bodyfont SGD iterates deviate from the full-batch gradient-flow path (right). For a gradient-step reference, match the update residual:
\[\color{WhartonBlue}-\frac{\Delta\theta_t}{\eta}-\nabla_\theta\mathcal L(\theta_t)
\;\approx\;\nabla_\theta\mathcal R(\theta_t;\Lambda),\]
where $\Delta\theta_t=\theta_{t+1}-\theta_t$.\\[0.45cm]
Pool steps to fit a shared regularizer, or let its parameters vary over time. A zero update recovers the endpoint case.\par}}
\block{49}{12}{42}{\heading{\icon{target}\quad Recovering known regularization}}
\block{49}{15.3}{42}{\textbf{Recover the penalty:} match known regularization parameters.\\[0.25cm]
\textbf{Recover the weights:} retrain with the estimated penalty to reproduce the learned weights.}
\block{49}{21}{42}{\subhead{Recovering elastic-net coefficients}}
\fig{49}{23.5}{42}{figures/elasticnet_recovery_mean_se.pdf}
\block{49}{37.2}{42}{{\captionfont Deep ReLU networks; $6\!\times\!6$ penalty grid, 10 datasets per cell. Mean $\pm$ SE; dashed line = exact recovery.}\\[0.55cm]
Recovery is accurate at intermediate strengths. Weak signals and nearly collinear gradients limit identification at the extremes.}
\block{49}{43.5}{42}{\subhead{Early stopping as a quadratic penalty}\\[0.4cm]
{\bodyfont The early-stopped OLS iterate solves a regularized problem~\cite{ali2019continuous}:}
{\bodyfont
\[\hat\theta_t=\arg\min_\theta\left\{\frac{1}{n}\|y-X\theta\|_2^2+\mathcal R_t(\theta)\right\},\]
\[\color{WhartonBlue}\mathcal R_t(\theta)=\theta^\top\Lambda^{(t)}\theta,\qquad
\Lambda^{(t)}=VS\big[(I-\eta S)^{-t}-I\big]^{-1}V^\top.\]}
{\captionfont Here $X^\top X/n=VSV^\top$, $\theta_0=0$, and $0<\eta s_i<1$.}}
\block{49}{54.2}{42}{{\bodyfont A symmetric $\Lambda$ has $p(p+1)/2$ unknowns but only $p$ equations per endpoint. For B, we pool endpoints with shared $X$ and $t$, varying the responses. Full-matrix identification requires diverse endpoints; a diagonal fit has only $p$ unknowns.}}
\fig{53.25}{58.5}{33.5}{figures/ols_composite.pdf}
\block{49}{78.7}{42}{{\captionfont
\textbf{A--B:} Theory (A) and recovery from 10 endpoints at the same stopping time (B).\\[0.15cm]
\textbf{C:} A diagonal fit from one endpoint is underparameterized, yet retraining reproduces the early-stopped weights: an equivalent regularizer, despite differing from the theoretical matrix.\\[0.15cm]
\textbf{D:} More endpoints improve matrix estimation, even when stopping times vary.\\[0.15cm]
\textbf{E:} The estimated scalar penalty decreases over training.\par}}
\block{94}{12}{42}{\heading{\icon{network}\quad Implicit regularization in deep networks}}
\block{94}{15.3}{42}{\subhead{Dropout rate and network architecture}}
\block{94}{17.5}{42}{{\bodyfont Dropout has data-dependent regularization effects~\cite{weiImplicitExplicitRegularization2020}. On MNIST, vary depth, width, and dropout rate; fit $\mathcal R(\theta,\lambda)=\lambda\|\theta\|_2^2$.}}
\fig{97.5}{21.1}{35}{figures/dropout_bias_ridge_panel.pdf}
\block{94}{43}{42}{{\captionfont Ten seeds per setting, including zero-dropout controls. Color shows gradient fit: lighter = more explained.}\\[0.6cm]
Narrow networks show an upturn at high dropout; wider networks have smaller estimates. Scalar ridge summarizes part of the effect.}
\block{94}{50}{42}{\subhead{Implicit gradient regularization}}
\block{94}{52.3}{42}{{\bodyfont Finite GD steps approximately follow a modified gradient flow with penalty $\mathcal R=\frac{\lambda}{p}\|\nabla\mathcal L\|^2$ and strength \mbox{$\lambda=\eta p/4$}~\cite{barrettImplicitGradientRegularization2022}. This penalizes steep parts of the trajectory, biasing learning toward flatter minima associated with better generalization. We recover $\lambda$ from trajectory deviations (right); stronger regularization accompanies smaller gradients and higher test accuracy (left, middle).}}
\fig{95.5}{59}{39}{figures/barrett_igr_figure2.pdf}
\block{94}{70}{42}{\subhead{Future work}\bodyfont
\begin{list}{\color{WhartonRed}\textbullet}{\setlength{\leftmargin}{1.1cm}\setlength{\labelsep}{0.35cm}\setlength{\itemsep}{0.2cm}\setlength{\topsep}{0.3cm}\setlength{\parsep}{0pt}}
\item \textbf{Diagnose training choices early}, before expensive full runs.
\item \textbf{Study complex architectures and optimizers} through their implicit penalties.
\item \textbf{Improve identification} with diverse endpoints and richer regularizer families.
\end{list}}
\block{94}{77}{42}{{\fontsize{28}{34}\selectfont\color{WhartonBlue}\bfseries References}}
\block{94}{78}{42}{\bibliography{references}}
\draw[WhartonBlue,line width=1.6pt] (4,-85) -- (136,-85);
\block{4}{86}{42}{{\fontsize{24}{30}\selectfont\icon[0.65cm]{file-text}\enspace\textbf{Paper:} \href{https://arxiv.org/abs/2605.05436}{arXiv:2605.05436}}}
\block{49}{86}{42}{{\fontsize{24}{30}\selectfont\icon[0.65cm]{book-open}\enspace\textbf{Blog:} \href{https://jrudoler.com/estimating-implicit-regularization/}{Project exposition}}}
\block{94}{86}{36}{{\fontsize{24}{30}\selectfont\icon[0.65cm]{code-xml}\enspace\textbf{Code:} \href{https://github.com/jrudoler/estimating-implicit-regularization}{GitHub repository}}}
\block{4}{88}{125}{{\fontsize{20}{25}\selectfont Supported in part by a Susquehanna PhD Fellowship awarded to Joseph H. Rudoler.}}
\node[anchor=north east,inner sep=0pt] at (136,-85.8) {\qrcode[height=2.7cm]{https://jrudoler.com/estimating-implicit-regularization/}};
\end{scope}\end{tikzpicture}
\end{document}
